\clearpage\SourceLeaf{418}%
\thispagestyle{empty}
\pdfbookmark[0]{Allgemeines Coordinaten-Verzeichniss.}{coordinaten-verzeichniss}
\begin{center}\Large ALLGEMEINES\par\medskip COORDINATEN-VERZEICHNISS\par\bigskip
\normalsize ZUSAMMENGETRAGEN\par\medskip AUS FOLGENDEN PARTIELLEN VERZEICHNISSEN\end{center}
\smallskip
\begin{list}{}{\setlength{\leftmargin}{2.5em}\setlength{\labelwidth}{2.3em}\setlength{\labelsep}{.2em}\setlength{\itemsep}{3pt}\setlength{\parsep}{0pt}}
\item[Nr. 1)] Generalverzeichniss von 1828. [Gradmessung 1821. 1822. 1823 und deren Fortsetzung bis Jever 1824. 1825 ausgeführt von C. F. \textsc{Gauss}.]
\item[2)] Eichsfeld 1828. [Messungen des Hauptmann \textsc{Müller} und des Lieutenant \textsc{Gauss}.]
\item[3)] Hildesheim 1828. [Messungen des Lieutenant \textsc{Hartmann}.]
\item[4)] Hildesheim 1829. [Messungen des Lieutenant \textsc{Hartmann}.]
\item[5)] Westphalen 1829. Messungen des Lieutenant \textsc{Gauss} und Lieutenant \textsc{Hartmann}.]
\item[6)] Westphalen 1830. [Messungen des Lieutenant \textsc{Hartmann}.]
\item[6*)] \textsc{Hartmanns} Messungen von 1831, so weit sie nicht durch 12 ersetzt werden.
\item[7)] Lüneburg. [Messungen des Hauptmann \textsc{Müller} im Jahre 1830 und des Lieutenant \textsc{Gauss} in den Jahren 1830 und 1831.]
\item[8)] Harz 1833. [Messungen des Lieutenant \textsc{Hartmann}.]
\item[9)] Mittelweser 1833. 1834. [Messungen des Hauptmann \textsc{Müller} und des Lieutenant \textsc{Gauss}.]
\item[10)] Oberweser (Göttingen) 1836. Messungen des Hauptmann \textsc{Müller}.]
\item[11)] Aller 1838. Messungen des Hauptmann \textsc{Müller}.]
\item[12)] Ostfriesland 1841. [Messungen des Hauptmann \textsc{Müller}.]
\item[13)] Bremen 1839. [Messungen des Hauptmann \textsc{Müller}.]
\item[14)] Bremen 1843. 1844. [Messungen des Lieutenant \textsc{Gauss}.]
\end{list}
\begin{center}1844. Dec. 13.\end{center}
